\part{参数高效微调与蒸馏}

\chapter{LoRA 与 QLoRA}
\label{sec:lora}

\begin{motivationbox}{每个任务都复制并更新完整模型过于昂贵}
全量微调需要为全部权重保存梯度和优化器状态，并为每个任务维护完整检查点。参数高效微调（Parameter-Efficient Fine-Tuning，PEFT）希望冻结大部分基础模型，只训练少量任务相关参数。LoRA 的关键假设是：预训练模型适配到某个下游任务时，所需权重更新常集中在较低维子空间，可以用低秩矩阵近似。
\end{motivationbox}

采用列向量约定，原线性层权重为
\[
W\in\mathbb R^{d_{\rm out}\times d_{\rm in}},
\qquad y=Wx.
\]
低秩适配（Low-Rank Adaptation，LoRA）冻结 $W$，只学习
\[
W'=W+\Delta W,
\qquad
\Delta W=\frac{\alpha}{r}BA,
\]
\[
A\in\mathbb R^{r\times d_{\rm in}},\qquad
B\in\mathbb R^{d_{\rm out}\times r},\qquad
r\ll\min(d_{\rm in},d_{\rm out}).
\]
这里 $r$ 是低秩维度，$\alpha$ 是缩放超参数，$\alpha/r$ 让更新幅度与 rank 的变化更容易解耦；$A$ 先把输入压到 $r$ 维，$B$ 再映射回输出维。不同实现的矩阵命名或转置约定可能相反，但参数量始终是
\[
r(d_{\rm in}+d_{\rm out}),
\]
而不是 $d_{\rm in}d_{\rm out}$。

\figLoRA

LoRA 常让一个低秩矩阵随机初始化、另一个初始化为零，使训练开始时 $\Delta W=0$，模型初始行为与基础模型完全一致，同时非零分支保证后续梯度能传播。若 $A$ 与 $B$ 都初始化为零，第一步两边梯度也可能同时为零，无法启动学习。

\begin{exercisebox}{4096 维方阵的 LoRA 参数量}
给定
\[
W\in\mathbb R^{4096\times4096},\qquad r=16.
\]
全量矩阵有
\[
4096^2=16{,}777{,}216
\]
个参数；LoRA 低秩分支有
\[
16(4096+4096)=131{,}072
\]
个参数，仅为完整矩阵的 $1/128$。这个比例描述该层的\emph{可训练权重}，不代表整次训练显存也精确下降 128 倍，因为基础权重、激活和临时缓冲仍然存在。
\end{exercisebox}

\codesection{对应实现：带缩放与零更新初始化的 LoRA 线性层}

\begin{lstlisting}
import math
import torch
import torch.nn as nn

class LoRALinear(nn.Module):
    def __init__(self, base: nn.Linear, rank=16, alpha=32.0):
        super().__init__()
        self.base = base
        self.rank = rank
        self.scale = alpha / rank

        # 冻结基础权重；训练时只更新 A、B。
        for parameter in self.base.parameters():
            parameter.requires_grad = False

        # PyTorch Linear 使用 y = x @ W.T。
        # A: [r, d_in]，B: [d_out, r]。
        self.lora_a = nn.Parameter(
            torch.empty(rank, base.in_features)
        )
        self.lora_b = nn.Parameter(
            torch.zeros(base.out_features, rank)
        )
        nn.init.kaiming_uniform_(self.lora_a, a=math.sqrt(5))

    def forward(self, x):
        base_output = self.base(x)

        # x @ A.T: [..., r]；再乘 B.T 回到 [..., d_out]。
        low_rank = (x @ self.lora_a.T) @ self.lora_b.T
        return base_output + self.scale * low_rank

    @torch.no_grad()
    def merged_weight(self):
        # 合并后恢复普通 Linear 权重，不再需要额外 LoRA 分支。
        delta = self.scale * (self.lora_b @ self.lora_a)
        return self.base.weight + delta
\end{lstlisting}

\section{LoRA 到底省什么？}

LoRA 最直接减少的是可训练参数、对应梯度和优化器状态；它也让每个任务只保存小型 adapter。基础模型前向与激活仍要计算，因此训练 FLOPs 和峰值显存不会按可训练参数比例等比下降。推理时可以把 $\Delta W$ 合并进 $W$，消除低秩分支的额外算子，但合并后的模型通常只是回到基础线性层速度，\textbf{LoRA 不自动让推理比原模型更快。}

Q、K、V、输出投影与 FFN 都可成为 LoRA 目标。Q/V 常被选中，是因为它们直接影响“如何查询”和“传递什么内容”，且参数预算较小；但最佳目标依任务而异，复杂领域适配可能需要 K、O 或 FFN。不能把“Q/V 最常见”误说成理论上唯一正确。

\begin{interviewbox}{为什么 LoRA 可以用低秩更新？}
\textbf{20 秒回答：}LoRA 假设下游适配不需要在完整权重空间任意移动，关键更新集中在较低维子空间，因此用 $BA$ 近似 $\Delta W$。这是一种经验归纳偏置，不是所有任务都保证成立；rank 太小会欠拟合，太大则逐渐失去参数效率。

\textbf{深入解释：}秩 $r$ 决定可表达更新子空间的维度，$\alpha/r$ 控制实际更新尺度。LoRA 主要节省可训练状态，而非消除基础模型计算。合并是把低秩乘积提前加到 $W$，使部署图恢复成普通线性层，并不改变合并后矩阵大小。

\textbf{面试官继续追问：}
\begin{enumerate}
  \item rank 越大一定越好吗？
  \item 为什么初始化时希望 $\Delta W=0$？
  \item $\alpha/r$ 的作用是什么？
  \item LoRA 与 Adapter 有什么区别？
  \item 哪些投影适合插 LoRA？
\end{enumerate}
\end{interviewbox}

\begin{wrongbox}
\textbf{错误说法：}LoRA 参数少，所以训练和推理都会按相同比例变快。

\textbf{更准确：}它显著减少梯度、优化器状态和任务检查点；基础模型矩阵乘法与激活仍存在。训练速度取决于激活反向、内核和数据流水线；推理合并后通常是不增加额外延迟，而不是天然快于基础模型。
\end{wrongbox}

\section{QLoRA}
\label{sec:qlora}

QLoRA 的核心组合是：\textbf{冻结量化后的基础模型，只训练高精度 LoRA adapter。}基础权重以低比特形式存储，前向时由量化计算内核使用或按需反量化；梯度不更新量化基础权重，而是流向 LoRA 参数。这样同时减少基础权重存储和可训练优化器状态，显著降低微调显存。

它与普通 INT4 推理量化不是同一件事。后者的目标是用低比特权重完成部署推理；QLoRA 的目标是在低比特冻结底座上进行参数高效训练。二者都涉及低比特基础权重，但训练图、梯度和 adapter 生命周期不同。

\begin{interviewbox}{基础权重量化后，为什么 LoRA 仍能训练？}
\textbf{20 秒回答：}基础权重被冻结，不需要对离散量化值做参数更新；前向仍能通过量化权重计算输出，损失梯度沿激活传播到旁路 LoRA 参数。真正更新的是通常以 BF16/FP16 保存的 $A,B$，不是 INT4 基础权重。

\textbf{深入解释：}量化误差改变了 LoRA 所适配的底座，因此量化格式、尺度和异常值处理仍会影响最终质量。QLoRA 节省的是底座存储和训练状态，不意味着所有激活都变成 4 bit，也不意味着训练过程与 INT4 推理完全相同。

\textbf{面试官继续追问：}
\begin{enumerate}
  \item QLoRA 真正训练哪些参数？
  \item 量化底座是否需要梯度？
  \item 为什么激活显存仍可能很大？
  \item QLoRA 与普通 INT4 inference quantization 有何区别？
\end{enumerate}
\end{interviewbox}

\chapter{知识蒸馏}
\label{sec:distillation}

\begin{motivationbox}{one-hot 标签没有表达教师知道的候选结构}
硬标签只告诉学生哪个类别被选中，不告诉它其他候选之间的相似性。大型教师模型的输出分布、生成回答和求解轨迹包含更丰富的结构。知识蒸馏用这些软信息训练更小学生模型，以容量和训练数据换取部署效率。
\end{motivationbox}

\section{Logits、回答与推理蒸馏}

设教师与学生 logits 分别为 $z_t,z_s$，温度为 $\tau>0$：
\[
p_t^{(\tau)}=\softmax(z_t/\tau),\qquad
p_s^{(\tau)}=\softmax(z_s/\tau).
\]
一个常见概念目标是
\[
\mathcal L=(1-\lambda)\mathcal L_{\rm hard}
+\lambda\tau^2D_{\rm KL}
\left(p_t^{(\tau)}\Vert p_s^{(\tau)}\right),
\]
其中 $\mathcal L_{\rm hard}$ 是真实标签交叉熵，$\lambda\in[0,1]$ 平衡硬标签与教师软分布，$\tau$ 把分布变平，使低概率候选之间的相对关系更可见；$\tau^2$ 常用于补偿温度导致的梯度尺度变化。

\textbf{Logits 蒸馏}匹配同一输入上的教师分布；\textbf{回答蒸馏}用大模型生成的数据训练小模型；\textbf{推理蒸馏}还使用经过验证的步骤或轨迹。大模型生成数据再训练小模型属于回答或序列级蒸馏，即使没有直接访问教师 logits，也仍是知识转移。

\begin{table}[htbp]
\centering
\small
\begin{tabularx}{0.94\textwidth}{lYY}
\toprule
方式 & 传递的信息 & 主要风险\\
\midrule
硬标签 & 最终正确类别 & 信息少，忽略候选关系\\
Logits 蒸馏 & 完整软类别分布 & 需要教师分数，词表与接口要对齐\\
回答蒸馏 & 教师生成的输入--输出对 & 继承风格与事实错误\\
推理蒸馏 & 答案及经过筛选的求解步骤 & 错误合理化、轨迹冗余、验证成本\\
\bottomrule
\end{tabularx}
\caption{不同蒸馏信号的比较。}
\end{table}

\begin{interviewbox}{教师 logits 比 one-hot 标签多提供了什么？}
\textbf{20 秒回答：}one-hot 只标出正确类，教师软分布还表达其他类别的相对相似度和不确定性。例如教师可能认为两个同义词都合理，只是一个概率更高；学生可以学习这种“暗知识”，而不是把所有非标签类别都同等视为错误。

\textbf{深入解释：}较高温度减小 logit 差，使长尾候选概率更可见，但温度过高会把有用排序抹平。蒸馏有效还依赖教师确实优于学生、训练输入覆盖部署分布，以及学生容量足以吸收目标；它不是无损压缩。

\textbf{面试官继续追问：}
\begin{enumerate}
  \item temperature 为什么用于 distillation？
  \item 回答蒸馏与 logits 蒸馏的接口要求有何不同？
  \item teacher 比 student 差时会怎样？
  \item reasoning distillation 有什么额外风险？
\end{enumerate}
\end{interviewbox}

推理蒸馏的更完整讨论见第\pageref{sec:reasoning-distillation}页。关键原则不变：教师轨迹必须经过答案验证和质量过滤，不能因为文本更长就视为更高质量监督。
